%%%%%%%%%%%%%%%%%%%%%%%%%%%%%%%%%%%%%%%%%%%%%%%%%%%%%%%%%%%%%%%%%%%%%%%
%
% Modelo de preenchimento de informações do cabeçalho
%
% Autor: Vitor G. P. Curtarelli
%
% Data: 10/03/2026
%
%%%%%%%%%%%%%%%%%%%%%%%%%%%%%%%%%%%%%%%%%%%%%%%%%%%%%%%%%%%%%%%%%%%%%%%

\titulo[Estimação do SSF na água com TT]{Desenvolvimentos na estimação do campo de velocidade do som na água utilizando tempo de trânsito}  % Título do relatório % Argumento em [] (opcional) é o título curto, a ser usado hos headers (vide começo da pág. 2)
%\subtitulo{}  % Sub-título do relatório, se houver
\autor{Vitor G. P. Curtarelli}  % Autor(es) do relatório
% Vários autores podem ser adicionados em um único comando \autor, ou o comando pode ser chamado várias vezes
%O número de autores no cabeçalho (Autor/Autores) detecta e atualiza automaticamente quando houver 2+ autores
\revisor{Stephan Paul, Lucas C. Lobato} % Revisor(es) do relatório
% Uso idêntico ao comando \autor
\entregavel{3.2.2}{Implementação dos algoritmos de TTT 2D/3D promissores identificados na Atividade 3.1}  % Número da atividade, como consta na Wiki do projeto
\resumo{\lipsum[2]}  % Um breve resumo do que o relatório trata
\comentario{Esse documento requer a leitura do relatório de revisão bibliográfica apresentado no Entregável 3.1.1, e \href{https://www.overleaf.com/read/zrwdjxsvnqvs\#e22068}{disponível no Overleaf}.}